\chapimg{fig/ch23.jpg}
\chapterlead{\subhd{After this chapter you can} compute the seven one-step effects and the arrow widths of the feedback layer, the nitrate cost of drainage, and the numbers behind each of the eight hidden questions.}
\chapter{Feedback loops and the hidden questions of the vadose zone}\label{ch:fb}
\markboth{Feedback loops}{Feedback loops}

\section{The seven effects}
For each control the engine moves one step on its list and re-evaluates Engine 1; the change is printed and the arrow width is $w_k=5+\mathrm{clip}(\mu_k;0,4)$ pixels, with $\mu_k$ from the table: temperature $\mu_1=8|\Delta Y|$; rain $\mu_2=|\Delta D|/12$; mulch $\mu_3=|\Delta E|/14$; trigger $\mu_4=|\Delta I|/8$; nitrate cost of drainage $\mu_5=|\Delta D|/10$; nitrogen $\mu_6=|\Delta N_L|/8$; consolidation $\mu_7=|\Delta D|/2$.

\exbox{Arrow widths}{One-step differences for a soil (illustration values: $\Delta Y=-0.11$ Mg ha$^{-1}$, $\Delta D=+9.0$ mm for the rain step, $\Delta E=-14$ mm, $\Delta I=-22$ mm, $\Delta N_L=+5.2$ kg ha$^{-1}$, $\Delta D=-0.6$ mm for consolidation) give the arrow widths of Table~\ref{tab:arrow}.}\sol  Example, row 4: $\mu_4=22/8=2.75$, so $w_4=5+2.75=7.75$ pixels. Row 3: $\mu_3=14/14=1.0$, $w_3=6.0$. Row 1: $\mu_1=8\times0.11=0.88$, $w_1=5.88$. A step that would give $\mu>4$ is clipped at width 9 pixels.\begin{table}[H]\centering\small\caption{Arrow widths of the seven feedbacks.}\label{tab:arrow}
\begin{tabular}{lrrrrr}\toprule
\hd{\#} & \hd{Control step} & \hd{Effect} & \hd{Value} & \hd{$\mu$} & \hd{Width (px)}\\\midrule
1 & Temperature $+0.5$ K & $\Delta Y$ & -0.11 & 0.880 & 5.88\\
2 & Rain $+0.1$ & $\Delta D$ & +9.00 & 0.750 & 5.75\\
3 & Stronger mulch & $\Delta E$ & -14.00 & 1.000 & 6.00\\
4 & Trigger one step drier & $\Delta I$ & -22.00 & 2.750 & 7.75\\
5 & Drainage (rain step) & $\Delta D$ & +9.00 & 0.900 & 5.90\\
6 & Nitrogen $+15$ per cent & $\Delta N_L$ & +5.20 & 0.650 & 5.65\\
7 & Consolidation one step & $\Delta D$ & -0.60 & 0.300 & 5.30\\\bottomrule\end{tabular}\end{table}


\exbox{The nitrate cost of drainage}{The study reports about $1.06$ kg N ha$^{-1}$ lost per mm of drainage: $\kappa_N=\dfrac{N_L^{(\text{rain step})}-N_L}{D^{(\text{rain step})}-D}$. With the rain step above ($\Delta D=+9.0$ mm, $\Delta N_L=+9.5$ kg N ha$^{-1}$), find $\kappa_N$ and the nitrate lost if the drainage is cut by 40 mm.}\sol  $\kappa_N=9.5/9.0=1.056$ kg N ha$^{-1}$ mm$^{-1}$, consistent with 1.06. Cutting drainage by 40 mm saves $40\times1.06=42.4$ kg N ha$^{-1}$, equal to a dressing of 42 kg of nitrogen that would otherwise be leached. The row is printed only if $|\Delta D|>0.5$ mm (otherwise ``no change''), because the ratio of two tiny differences is unstable.


\section{Hidden question 1: wetting and drying}
The wetting curve has twice the $\alpha$ of the drying curve. A sensor reads a suction from the water content; the same water content therefore corresponds to different suctions on the two branches.

\exbox{The same water content on two branches}{The loam ($\alpha_d=0.036$, $n=1.56$) is at the trigger on the drying curve, $h_d=300$ cm. Find the suction that the same water content has on the wetting curve ($\alpha_w=2\alpha_d$).}\sol  $\Se=\Se(300)=0.2615$ on the drying curve. On the wetting branch, $h_w=\alpha_w^{-1}(\Se^{-1/m}-1)^{1/n}$; since $\alpha_w=2\alpha_d$, $h_w=h_d/2=150$ cm exactly (check: $150.0$ cm). A sensor that has just been wetted thus reads 150 cm where a drying soil at the same water content reads 300 cm: the trigger set in suction fires at twice the water deficit after wetting as after drying. Hysteresis makes the same reading mean different amounts of water.


\section{Hidden question 2: where does water still move}
The capillary conductivity is $K_s K_r(h)$, and a film term proportional to $h^{-1.5}$ and anchored at $10^{-7}$ cm d$^{-1}$ at pF 4.2 ($h=15\,849$ cm) is added.

\exbox{Capillary and film conductivity}{For the loam with $K_s=24.96$ cm d$^{-1}$ compute the capillary conductivity, the film term and their sum at suctions 300, 3000, 15,849, $10^5$ and 316,228 cm (pF 5.5), and find the suction at which the film term overtakes the capillary term.}\sol  At pF 4.2 the film term is $10^{-7}$ cm d$^{-1}$ by construction. At pF 5.5: $10^{-7}(316228/15849)^{-1.5}=10^{-7}\times0.01122=1.12\times10^{-9}$ cm d$^{-1}$. The capillary term is $K_sK_r(h)$ with $K_r$ from the Mualem equation.\begin{table}[H]\centering\small\caption{Conductivity (cm d$^{-1}$) of the loam, capillary and film terms.}\label{tab:film}
\begin{tabular}{lrrr}\toprule
\hd{$h$ (cm)} & \hd{Capillary} & \hd{Film} & \hd{Sum}\\\midrule
300 & 9.50e-04 & 3.84e-05 & 9.88e-04\\
3{,}000 & 3.92e-07 & 1.21e-06 & 1.61e-06\\
15{,}849 & 1.37e-09 & 1.00e-07 & 1.01e-07\\
100{,}000 & 2.61e-12 & 6.31e-09 & 6.31e-09\\
316{,}228 & 5.20e-14 & 1.12e-09 & 1.12e-09\\\bottomrule\end{tabular}\end{table}
The two terms are equal at $h\approx1{,}653$ cm; beyond it, the film dominates. Where the soil is drier than that the conductivity falls only as $h^{-1.5}$ (a factor of 31.6 per decade of suction) instead of by orders of magnitude, so some water still moves in air-dry soil. The film term is an \emph{illustrative} form, not calibrated.


\section{Hidden question 5: where does the nitrate travel}
The panel shows the pores that empty at less than 31.6 cm (the \emph{wide} pores, radius above $47\ \mu$m) and the share of pore volume and of conductivity that they carry: share of volume $1-\Se(31.6)$; share of conductivity $1-K_r(31.6)$.

\exbox{Wide pores and conductivity}{Compute these two shares for the ten soils.}\sol  For the loam: $\Se(31.6)=0.751$ and $K_r(31.6)=0.032$; the wide pores hold 25 per cent of the pore volume and carry 96.8 per cent of the conductivity.\begin{table}[H]\centering\small\caption{Share of pore volume and of conductivity in pores emptied below 31.6 cm.}\label{tab:wide}
\begin{tabular}{lrrr}\toprule
\hd{ID} & \hd{Soil} & \hd{Volume share} & \hd{Conductivity share}\\\midrule
S01 & Coarse loam & 0.359 & 0.9979\\
S02 & Coarse sand & 0.363 & 0.9874\\
S03 & Sand & 0.923 & 1.0000\\
S04 & Loamy sand & 0.830 & 0.9998\\
S05 & Sandy loam & 0.574 & 0.9957\\
S06 & Loam & 0.249 & 0.9677\\
S07 & Silt loam & 0.115 & 0.9330\\
S08 & Sandy clay loam & 0.335 & 0.9914\\
S09 & Clay loam & 0.093 & 0.9514\\
S10 & Silty clay loam & 0.040 & 0.9321\\\bottomrule\end{tabular}\end{table}
\textbf{Reading.} Where the wide pores carry nearly all of the conductivity, solute moves fast through a small part of the pore volume, so a nitrate pulse reaches the drain earlier than the mean water content suggests.


\section{Hidden question 4: the water table}
The water content above a water table at height $z$ is $\Se=[1+(\alpha z)^n]^{-m}$ (equilibrium, suction $h=z$); the quasi-saturated height is where $\Se<0.9$, that is $z_{90}=\alpha^{-1}(0.9^{-1/m}-1)^{1/n}$.

\exbox{How far the water table reaches into the root zone}{Find $\Se$ at 5, 25, 60 and 150 cm above the water table and the height $z_{90}$ for the ten soils.}\sol  For the loam ($\alpha=0.036$, $n=1.56$, $m=0.359$): $0.9^{-1/m}=0.9^{-2.786}=1.341$, so $z_{90}=(1.341-1)^{0.641}/0.036=(0.341)^{0.641}/0.036=0.5013/0.036=13.9$ cm. Table~\ref{tab:wt} lists the ten soils.\begin{table}[H]\centering\small\caption{Water content above a water table and quasi-saturated height.}\label{tab:wt}
\begin{tabular}{lrrrrrr}\toprule
\hd{ID} & \hd{Soil} & \hd{$\Se$(5)} & \hd{$\Se$(25)} & \hd{$\Se$(60)} & \hd{$\Se$(150)} & \hd{$z_{90}$ (cm)}\\\midrule
S01 & Coarse loam & 0.921 & 0.685 & 0.526 & 0.387 & 6.2\\
S02 & Coarse sand & 0.958 & 0.699 & 0.470 & 0.283 & 9.2\\
S03 & Sand & 0.802 & 0.113 & 0.026 & 0.006 & 3.7\\
S04 & Loamy sand & 0.850 & 0.226 & 0.076 & 0.024 & 4.0\\
S05 & Sandy loam & 0.934 & 0.504 & 0.255 & 0.115 & 6.4\\
S06 & Loam & 0.976 & 0.802 & 0.591 & 0.379 & 13.9\\
S07 & Silt loam & 0.989 & 0.911 & 0.786 & 0.603 & 27.8\\
S08 & Sandy clay loam & 0.952 & 0.718 & 0.520 & 0.347 & 8.9\\
S09 & Clay loam & 0.989 & 0.927 & 0.831 & 0.685 & 33.8\\
S10 & Silty clay loam & 0.995 & 0.969 & 0.923 & 0.834 & 79.7\\\bottomrule\end{tabular}\end{table}


\section{Hidden questions 6 and 7: salinity and the uptake curve}
At steady state the salt entering equals the salt leaving: $I\,EC_{\rm irr}=D\,EC_{\rm dr}$. The root-zone salinity is taken as a quarter of the sum, $EC_{\rm rz}=(EC_{\rm irr}+EC_{\rm dr})/4$, and wheat loses $7.1$ per cent of yield per dS m$^{-1}$ above a threshold of 6 dS m$^{-1}$.

\exbox{How salty can the irrigation water be}{For an irrigation of $I=150$ mm and drainage of $60$, $30$ and $15$ mm, find $EC_{\rm dr}$, $EC_{\rm rz}$ and the relative yield for irrigation water of 2, 4 and 6 dS m$^{-1}$. Then find the largest irrigation-water salinity that keeps the yield at 100 per cent.}\sol  $EC_{\rm dr}=EC_{\rm irr}I/D$. For $I=150$, $D=60$, $EC_{\rm irr}=2$: $EC_{\rm dr}=5$, $EC_{\rm rz}=(2+5)/4=1.75$ dS m$^{-1}$, below 6, so yield 100 per cent. For $D=15$ and $EC_{\rm irr}=6$: $EC_{\rm dr}=60$, $EC_{\rm rz}=16.5$, loss $0.071(16.5-6)=0.745$, yield 25.5 per cent. See Table~\ref{tab:salt}.\begin{table}[H]\centering\small\caption{Salt balance and wheat yield.}\label{tab:salt}
\begin{tabular}{lrrrrr}\toprule
\hd{$I$} & \hd{$D$} & \hd{$EC_{\rm irr}$} & \hd{$EC_{\rm dr}$} & \hd{$EC_{\rm rz}$} & \hd{Yield (\%)}\\\midrule
150 & 60 & 2 & 5.0 & 1.75 & 100.0\\
150 & 60 & 4 & 10.0 & 3.50 & 100.0\\
150 & 60 & 6 & 15.0 & 5.25 & 100.0\\
150 & 30 & 2 & 10.0 & 3.00 & 100.0\\
150 & 30 & 4 & 20.0 & 6.00 & 100.0\\
150 & 30 & 6 & 30.0 & 9.00 & 78.7\\
150 & 15 & 2 & 20.0 & 5.50 & 100.0\\
150 & 15 & 4 & 40.0 & 11.00 & 64.5\\
150 & 15 & 6 & 60.0 & 16.50 & 25.5\\\bottomrule\end{tabular}\end{table}
\textbf{Threshold.} $EC_{\rm rz}=EC_{\rm irr}(1+I/D)/4\le6$ gives $EC_{\rm irr}\le24/(1+I/D)$: $6.86$ dS m$^{-1}$ for $D=60$ ($I/D=2.5$), $4.0$ for $D=30$ and $2.2$ for $D=15$. Less drainage means less leaching and a lower tolerable salinity: the same irrigation water that is safe in a free-draining soil damages wheat in a tight soil. The steady state, the one-quarter rule and the threshold are stated simple forms ``for teaching and comparison'', not calibrated to a field.


\exbox{Reduction of root uptake with suction}{The uptake curve has breakpoints at 10, 25, 500 and 8,000 cm (a Feddes-type reduction function; here assumed linear in $h$ between breakpoints): uptake is zero at $h\le10$ cm (anaerobic), rises linearly to 1 at 25 cm, stays at 1 until 500 cm and falls linearly to 0 at 8,000 cm. Find the relative uptake at $h=15,\ 300,\ 1000,\ 3000$ and 6000 cm.}\sol  At 15 cm: $(15-10)/(25-10)=0.333$. At 300 cm: 1. At 1000 cm: $1-(1000-500)/7500=0.933$. At 3000: $1-2500/7500=0.667$. At 6000: $1-5500/7500=0.267$. A trigger of $-1000$ cm therefore costs a crop only about 7 per cent of the potential uptake at the moment of the trigger, but the average over the dry-down is lower. In the loam $-1000$ cm is a depletion of 0.76, beyond the $f^*=0.45$.


\section{Hidden question 8: keeping up with the sun}
The topsoil delivers water against the demand until stage 2 begins; mulch reduces the demand and lengthens stage 1 (Chapter~\ref{ch:mulch}). The panel's answer is the stage-1 duration $t_1=Q_1/e_0$.
\exbox{Stage 1 with and without mulch}{For the sand (S03, $Q_1=12.7$ mm) and the silty clay loam (S10, $Q_1=4.6$ mm) find how long the topsoil keeps up with a demand of 4.2 mm d$^{-1}$ bare, and with a mulch $\varphi=0.3$.}\sol  Bare: sand $12.7/4.2=3.0$ d; silty clay loam $4.6/4.2=1.1$ d. With mulch $e_0=4.2\times0.3=1.26$: sand $12.7/1.26=10.1$ d; silty clay loam $4.6/1.26=3.7$ d. Mulch triples the time that the surface keeps pace in both soils; the sand has the longer reserve.

